\documentclass[11pt]{article}
\usepackage[margin=1in]{geometry}
\usepackage{amsmath,amssymb,mathtools}
\usepackage{fontspec,unicode-math}
\setmainfont{Latin Modern Roman}
\setsansfont{Latin Modern Sans}
\setmathfont{Latin Modern Math}
\usepackage{xurl,hyperref}
\hypersetup{colorlinks=true,urlcolor=blue,linkcolor=blue}
\newcommand{\sref}[2]{\hyperlink{source:#1:#2}{[S#2]}}
\newcommand{\cataloguescope}[1]{\paragraph{Recorded mathematical question.}#1}
\begin{document}
\section{Finite-dimensional distance laws of the Brownian map}\label{statement}

\subsection{Definitions and mathematical statement}

\paragraph{Shared notation.}$\mathbb N=\{1,2,\ldots\}$, and $\mathbb Z,\mathbb Q,\mathbb R,\mathbb C$ denote the integers, rationals, reals and complex numbers. Any other conventions are specified locally.


Let $(\mathcal M,d,\mu)$ be a Brownian map with its volume measure normalized to a probability measure. This random measured compact metric space is obtained from scaling limits of uniform finite planar maps equipped with graph distance multiplied by the conventional $n^{-1/4}$ factor and uniform vertex probability measure. Conditional on $(\mathcal M,d,\mu)$, sample $X_1,\ldots,X_p$ independently from $\mu$, and form $\mathsf D_p=(d(X_i,X_j))_{1\le i,j\le p}$.
\paragraph{Statement recorded in the supplied source.}
For every integer $p\ge4$, determine the joint probability distribution of all entries of $\mathsf D_p$. Any choice of overall Brownian-map distance normalization should be stated explicitly in a resulting formula. \sref{statement}{2}
\subsection{Short English statement}
Find the joint distance law for four or more independent volume-sampled points of the Brownian map.
\subsection{Sources}
\begin{description}
\item[\hypertarget{source:statement:1}{S1}] ULAM.AI and the UnsolvedMath Contributors, \emph{UnsolvedMath: A Curated Collection of Open Mathematics Problems}, catalogue entry. CC BY 4.0. \url{https://huggingface.co/datasets/ulamai/UnsolvedMath}.
\item[\hypertarget{source:statement:2}{S2}] Itai Benjamini, \emph{Random Planar Metrics}; original author manuscript preserved in Arquivo.pt capture of 31 December 2020, Section 3.1, extension (2), p. 5. \url{https://arquivo.pt/noFrame/replay/20201231041544id_/http://www.wisdom.weizmann.ac.il/~itai/randomplanar2.pdf}.
\end{description}
\label{end:statement}
\end{document}
